% This file is part of the modular source for Stratified Relational Workspace Theory.
% !TeX root = ../Stratified_Relational_Workspace_Theory.tex
\theorypart{V}{Gauge Invariance and Abstract Identity}

Exact compatibility is a coordinate-free notion. Approximate spectra also
depend on the Hilbert metrics used to measure discrepancy. Gauge theory must
therefore distinguish a change of coordinates from a physical change of
metric.

\topic{16}{Local Coordinate Transformations}

For each \(i\in V\), let
\[
g_i\in GL(C_i),
\qquad
g=(g_i)_{i\in V}.
\]
Define the vertex and edge direct sums
\[
G_V=\bigoplus_{i\in V}g_i,
\qquad
G_E=\bigoplus_{e\in E}g_{t(e)}.
\]
The transformed transfers and local operation representations are
\[
H_e^g
=
 g_{t(e)}H_eg_{s(e)}^{-1},
\]
\[
\rho_i^g(a)
=
 g_i\rho_i(a)g_i^{-1}.
\]
The fundamental covariance identity is
\[
\boxed{
D_{H^g}G_V=G_ED_H.
}
\]
It follows directly by evaluating the \(e\)-component:
\[
\sqrt{w_e}
\left(
 g_{t(e)}c_{t(e)}
-
 g_{t(e)}H_ec_{s(e)}
\right).
\]

\topic{17}{Gauge Invariance Theorem}

\begin{theorem}[Exact gauge covariance]\label{thm:gauge-covariance}
The map
\[
G_V:\Gamma(H)\longrightarrow\Gamma(H^g)
\]
is an isomorphism. Therefore,
\[
\boxed{
\Gamma(H^g)=G_V\Gamma(H),
\qquad
\dim\Gamma(H^g)=\dim\Gamma(H).
}
\]
If the core transports are invertible and \(o\) is a root, then
\[
H_\gamma^g=g_oH_\gamma g_o^{-1}
\]
for every loop \(\gamma:o\to o\). Hence the holonomy-fixed core transforms by
\[
C_*^g=g_oC_*.
\]
\end{theorem}

\begin{proof}
The covariance identity implies
\[
D_{H^g}(G_Vc)=G_ED_Hc.
\]
Since \(G_E\) is invertible,
\[
D_Hc=0
\quad\Longleftrightarrow\quad
D_{H^g}(G_Vc)=0.
\]
The holonomy formula follows by cancellation of adjacent factors
\(g_i^{-1}g_i\) along a path.
\end{proof}

\begin{theorem}[Unitary gauge covariance of the Laplacian]\label{thm:unitary-gauge}
If every \(g_i\) is unitary, then \(G_V\) and \(G_E\) are unitary and
\[
\boxed{
L_{H^g}=G_VL_HG_V^{-1}.
}
\]
Consequently,
\[
\spec(L_{H^g})=\spec(L_H),
\qquad
P_I(L_{H^g})=G_VP_I(L_H)G_V^{-1}
\]
for every spectral interval \(I\).
\end{theorem}

\begin{proof}
From
\[
D_{H^g}=G_ED_HG_V^{-1}
\]
and unitarity,
\[
D_{H^g}^*D_{H^g}
=
G_VD_H^*G_E^{-1}G_ED_HG_V^{-1}
=
G_VL_HG_V^{-1}.
\]
The spectral statements follow from unitary similarity.
\end{proof}

\subtopic{17.1 Nonunitary gauges and metric distortion}

For a general invertible gauge, define
\[
m_V=\sigma_{\min}(G_V),
\qquad
M_V=\|G_V\|,
\]
\[
m_E=\sigma_{\min}(G_E),
\qquad
M_E=\|G_E\|.
\]
The covariance identity yields
\[
m_E^2\mathcal E_H(c)
\leq
\mathcal E_{H^g}(G_Vc)
\leq
M_E^2\mathcal E_H(c),
\]
and
\[
m_V^2\|c\|^2
\leq
\|G_Vc\|^2
\leq
M_V^2\|c\|^2.
\]
Therefore the Rayleigh quotients obey
\[
\boxed{
\frac{m_E^2}{M_V^2}
\frac{\mathcal E_H(c)}{\|c\|^2}
\leq
\frac{\mathcal E_{H^g}(G_Vc)}{\|G_Vc\|^2}
\leq
\frac{M_E^2}{m_V^2}
\frac{\mathcal E_H(c)}{\|c\|^2}.
}
\]
By the min--max principle,
\[
\boxed{
\frac{m_E^2}{M_V^2}\lambda_k(L_H)
\leq
\lambda_k(L_{H^g})
\leq
\frac{M_E^2}{m_V^2}\lambda_k(L_H).
}
\]
Exact zero modes are unchanged, but fixed-metric approximate eigenvalues can be
strongly distorted by an ill-conditioned coordinate transformation.

\subtopic{17.2 Metric-covariant formulation}

A pure change of coordinates should transform the metric as well as the
matrices. On each transformed space, define
\[
\langle u,v\rangle_i^g
=
\left\langle g_i^{-1}u,g_i^{-1}v\right\rangle_i.
\]
Use the analogous pullback metric on every edge copy. Then \(G_V\) and
\(G_E\) are unitary between the original and transformed metric spaces.
Consequently, the complete transfer-Laplacian spectrum is exactly preserved.

This gives two distinct interpretations:
\begin{itemize}
\item Transforming the matrices and the metrics gives the same approximate
workspace in new coordinates.
\item Transforming the matrices while holding the metrics fixed gives a
potentially different weighting of discrepancy.
\end{itemize}
The distinction prevents arbitrary coordinate scaling from manufacturing or
erasing apparent low-energy workspace directions.

\topic{18}{Static Workspace Equivalence}

\begin{definition}[Exact static equivalence]\label{def:static-equivalence}
Two static realizations are exactly gauge equivalent,
\[
\mathfrak R\sim_{\mathrm{stat}}\mathfrak R',
\]
when there exists an invertible family \(g_i\) such that
\[
H_e'=g_{t(e)}H_eg_{s(e)}^{-1},
\qquad
\rho_i'(a)=g_i\rho_i(a)g_i^{-1}.
\]
The exact static abstract workspace is
\[
\boxed{
\Workspace_{\mathrm{stat}}
=
[\mathfrak R]_{\sim_{\mathrm{stat}}}.
}
\]
\end{definition}

For approximate theory, the realization also contains the vertex and edge
metrics. Two metric realizations are equivalent only when the gauge transports
those metrics, or when the gauge is unitary relative to the fixed metrics.
Accordingly:
\[
\boxed{
\begin{aligned}
\text{exact workspace identity}
&\text{ is invariant under }GL\text{-gauge},\\
\text{fixed-metric spectral geometry}
&\text{ is invariant under unitary gauge},\\
\text{metric-covariant spectral geometry}
&\text{ is invariant under transported }GL\text{-gauge}.
\end{aligned}
}
\]

\clearpage
\subtopic{18.1 What the approximate theory establishes}

Parts~IV--V now prove:
\[
\begin{aligned}
&\ker L_H=\Gamma(H),\\
&\text{low energy plus a gap implies proximity to an exact core},\\
&\text{a separated workspace cluster is stable under transfer noise},\\
&\text{approximately intertwined operations nearly preserve that cluster},\\
&\text{core leakage and residual reentry admit explicit pathwise bounds},\\
&\text{exact identity is gauge invariant while approximate geometry is metric aware}.
\end{aligned}
\]
These are structural and perturbative results. They still do not prove that a
workspace variable is causally used, predictively sufficient, irreducibly
distributed, or temporally self-maintaining. Those questions begin in
Part~VI.

\subtopic{18.2 Finite computation protocol}

For a declared graph, set of core dimensions, weights, and transfer maps, the
static analysis can be executed in the following order:
\begin{enumerate}
\item Assemble the sparse block operator \(D_H\) from the edge equations
\[
\sqrt{w_e}\bigl(c_{t(e)}-H_ec_{s(e)}\bigr).
\]
\item Compute a singular-value decomposition
\[
D_H=U\Sigma V^*,
\qquad
\lambda_k(L_H)=\sigma_k(D_H)^2.
\]
Use a declared numerical tolerance to distinguish computational zero modes
from positive modes.
\item Select a candidate rank \(d\) only at a reproducible spectral boundary
with \(g_d>0\). A large observed gap is evidence of stability, not evidence of
causal relevance.
\item Quantify transfer uncertainty by \(\delta_H\) and verify whether
\[
2\|D_H\|\delta_H+\delta_H^2<\frac{g_d}{2}.
\]
If this gate fails, the extracted projector has no certified perturbative
identity at the stated error scale.
\item For every proposed common operation, compute \(\beta(a)\) and the
operation-leakage bound in Theorem~\ref{thm:approx-operation-closure}.
\item For every long transfer chain used in an interpretation, report the
pathwise leakage certificate from Theorem~\ref{thm:residual-duhamel} rather
than extrapolating from single-edge defects.
\end{enumerate}

A minimal structural report is therefore
\[
\boxed{
\left(
 d,\lambda_d,\lambda_{d+1},g_d,
 \delta_H,
 \|P_d-\widetilde P_d\|_{\mathrm{bound}},
 \beta(a),
 \Lambda_\gamma
\right),
}
\]
where \(\Lambda_\gamma\) denotes the right-hand side of the pathwise leakage
bound. This structural report is intentionally incomplete as a workspace
certificate. Parts~VI--IX add the predictive, causal, null, and irreducibility
gates required to interpret the low-energy subspace operationally.
